% Procedencia primaria (lectura de 18-09-2026):
% XROOT = output/EXTREMA_MEDIA_RAZON_NARRACION_Y_REVISION_20260916/ARTICULO/ES/source
% XROOT/libro/01_acoplamiento.tex:51-147,727-749.
% XROOT/sections/k_direccion_dimensional.tex:54-155.
% XROOT/nuclear/12.tex:9-88.
% VIROOT = output/AMPLIACION_FORMAL_LEAN_SERIE_20260916/delivery/USB_ES_EN_STAGE10_20260917/ES/PAQUETES/06_PARTICULAS_Y_MASAS_ES/documentacion_original/payload/source_es
% VIROOT/sections/05_registro_operador_masa.tex:11-119,193-294.
% VIROOT/sections/02_particula_persistencia.tex:210-235.
% XROOT/nuclear/07.tex:5-25,74-154: Hamiltoniano de transporte eliptico.
% XROOT/nuclear/11.tex:878-979: elevacion cotangente y accion discreta.
% Estatutos: cadena K->alpha->accion->masa = RESULTADO_RECUPERADO;
% carta 1+1+10 de eventos y criterio de entrelazamiento = FORMALIZACION_NUEVA;
% congruencia, espectro generalizado y control racional = pruebas completas aqui.

\section{Positive congruences, action readers and Hamiltonian dynamics}
\label{sec:k-mass-bridge}

The relation between the dodecaphasic record and the General Mass Law has
an effective composition. APP produces the two sheets with their residues and
quotients; TRIT orients the transition and determines its regime; TPK composes
selection, transport and updating while preserving route, carry, boundary,
survival and memory. The regional readouts and the twelve-position record
arise jointly from the enriched state and its compatible
prolongation. Numerical closure and the material reading
share this origin and retain their specific operations.

With the phase origin, orientation and base one thousand declared, the record is
\[
 K=(234,543,140,729,659,824,621,58,914,794,146,601).
\]
The positional reader uses
\begin{equation}
 \kappa_{\rm ag}(K)=\frac{\sum_{j=1}^{12}K_j1000^{12-j}}{1000^{12}-1},
 \qquad
 \alpha=\pi_{\rm HMT}+e_{\rm HMT}-\varphi_{\rm HMT}-4-
              \kappa_{\rm ag}(K).
 \label{eq:mass-k-alpha}
\end{equation}
The values on the right are internal readouts of the same coinductive
process. The formula expresses the compatible infinite-closure reader,
with the carry boundary that defines it.

Action receives these readouts through the readers
\begin{align}
 H_5={}&\frac12\sqrt{1000\alpha/\varphi}-\alpha+
       \frac9{16}\alpha^2-\frac59\alpha^3+
       \frac7{48}\alpha^4-\frac1{54}\alpha^5,\notag\\
 D_A={}&\exp(-100\pi\alpha/9),\qquad
 \mathcal C_\pi=169\alpha^6+\frac{D_A\alpha^7}{1-D_A\alpha},\notag\\
 \eta_{\rm ret}={}&H_5-\frac{90}{\pi}\mathcal C_\pi,
 \qquad R_{\rm act}=\frac{H_5}{\eta_{\rm ret}}>0.
 \label{eq:mass-k-action}
\end{align}
Here and in the following formulae, \(\pi,\varphi,e\) abbreviate the
HMT readouts already constructed. The local determinantal norm
\(\Sigma_H=\varphi\alpha^{16}\) fixes the decade \(-34\), and the
section \(\hbar_{\rm ret}=\eta_{\rm ret}10^{-34}\mathcal U_S\)
belongs to the declared action line. The ratio \(R_{\rm act}\)
preserves the relation between the two action sections when their common
dimensional basis changes.

The central electron route provides, through its displacement and
autocorrelation readers, the triple \((80,54,6)\). With
\[
 \Delta_4=\frac{\pi-e\log\pi}{270}\left(1+\frac\pi{729}\right),
 \qquad
 b_e^{\rm alg}=\frac{\sqrt3}{4}
       \bigl[\exp(\varphi/\pi^2)-22\alpha^3\bigr](1+15\Delta_4),
\]
its material readout contains
\begin{equation}
 \kappa_e=80-54\alpha+6\Delta_4,
 \qquad \varepsilon_e^\beta=b_e^{\rm alg}R_{\rm act}^{\kappa_e}.
 \label{eq:mass-k-electron}
\end{equation}
The composition \eqref{eq:mass-k-alpha}--\eqref{eq:mass-k-electron}
identifies the specific role of \(K\): its closure reading
contributes to \(\alpha\), and this output subsequently acts in the action
reader and the return exponent. The central selection of the electron
belongs to the involution of the APP chart, whose unique fixed point is
\((5,5)\). The centrality of that cell and its scalar evaluation are
two results linked by the route and its record.

\subsection{Reversible coordinates of the material event record}

The material record assigns to a route \(\gamma\) twelve oriented
events \(\mathbf e(\gamma)=(e_0,\ldots,e_{11})\), one for each
nonadic window of the one-hundred-and-eight-position calendar. The same
temporal order allows these events to be realized in the twelve-coordinate
carrier used by \(K\). This identification requires preserving
the origin and orientation of the calendar.

Let \(\mathbf1\) be the uniform vector,
\(P_{11}=I-\mathbf1\mathbf1^*/12\), and let \(P_3\) be the
three-dimensional projector of the declared incidence chart. The geometric
construction of the record determines
\[
 u_K=P_3P_{11}K,\qquad
 \|u_K\|^2=\frac{6638585+2275584\sqrt5}{20}>0,
 \qquad u_K\perp\mathbf1,
\]
and the transverse projector
\[
 P_{10}(K)=P_{11}-\frac{u_Ku_K^*}{\|u_K\|^2}.
\]

\begin{proposition}[Material decomposition relative to the direction of \(K\)]
For every event vector \(\mathbf e\), the map
\begin{equation}
 \mathbf e\longmapsto
 \left(s_C=\mathbf1^*\mathbf e,\quad
       a_K=\frac{u_K^*\mathbf e}{\|u_K\|^2},\quad
       v_K=P_{10}(K)\mathbf e\right)
 \label{eq:mass-k-event-chart}
\end{equation}
is invertible onto its image and satisfies
\begin{equation}
 \mathbf e=\frac{s_C}{12}\mathbf1+a_Ku_K+v_K,
 \qquad
 \|\mathbf e\|^2=\frac{|s_C|^2}{12}
          +|a_K|^2\|u_K\|^2+\|v_K\|^2.
 \label{eq:mass-k-event-inverse}
\end{equation}
\end{proposition}
\begin{proof}
The projectors \(\mathbf1\mathbf1^*/12\),
\(u_Ku_K^*/\|u_K\|^2\) and \(P_{10}(K)\) are pairwise
orthogonal and sum to the identity. Applying this equality to \(\mathbf e\)
proves the inversion. Orthogonality of the three summands gives the
norm identity.
\end{proof}

This chart preserves the complete event vector. It therefore also preserves
the functionals that depend on its signs, sums and
correlations, once reconstruction
\eqref{eq:mass-k-event-inverse} has been performed. The record of predecessors of the nines,
the action count and the tower coordinates remain
in the enriched trajectory: the preceding chart modifies only
the coordinates of the twelve events. The character is evaluated on the same
recovered signature. The construction thus provides a geometric
representation of the information used by mass, with a direction
distinguished by \(K\) and ten transverse components.

The amount of information retained by each readout can be established
exactly. In the algebraic domain of records, \(K'=K+\mathbf1\)
satisfies
\[
 u_{K'}=u_K,\qquad P_{10}(K')=P_{10}(K),\qquad
 \kappa_{\rm ag}(K')-\kappa_{\rm ag}(K)=\frac1{999}.
\]
The last identity follows by summing the geometric progression of the
twelve positional weights. If the regional readouts are held
fixed, the reading of \(\alpha\) changes by \(-1/999\).
This algebraic check determines the precise scope of the direction:
its projector preserves a polarization, whereas scalar closure
also uses the uniform and positional information of the record.
The admissibility of \(K'\) as another terminal history requires its
own APP--TRIT--TPK construction.

\subsection{Positive congruence and spectrum of the material operator}

For a multisection \(M\), let \(\mathscr H_M\) be the finite-dimensional space
with an orthonormal basis labelled by its routes. The return reader
\(r\in\{D,\Omega\}\) produces a self-adjoint operator
\(B_M^r e_\gamma=\kappa_\gamma^r e_\gamma\). The refined character,
evaluated on differences relative to the electron, gives the basal operator
\[
 M_M^{(0)}e_\gamma=
 m_e^{(0)}\chi_{\rm ref}(\sigma_\gamma-\sigma_e,
                  \eta_\gamma-\eta_e)e_\gamma.
\]
The tower coefficients and their convergence belong to the reader of each
route. Its beta readout is the congruence
\begin{equation}
 D_M^r=R_{\rm act}^{B_M^r/2},\qquad
 M_M^{\beta,r}=D_M^rM_M^{(0)}D_M^r.
 \label{eq:mass-k-congruence}
\end{equation}

\begin{proposition}[Exact invariants of material congruence]
The operator \(D_M^r\) is positive and invertible. The congruence
\eqref{eq:mass-k-congruence} preserves rank and inertia; in particular,
it preserves positivity. If \(M_M^{(0)}\) and \(B_M^r\) are diagonal
in the route basis, then
\[
 m_\gamma^\beta=m_\gamma^{(0)}R_{\rm act}^{\kappa_\gamma^r}.
\]
For two routes in the same dimensional chart, one obtains
\[
 \frac{m_Y^\beta}{m_X^\beta}=
 \chi_{\rm ref}(\sigma_Y-\sigma_X,\eta_Y-\eta_X)
 R_{\rm act}^{\kappa_Y^r-\kappa_X^r}.
\]
\end{proposition}
\begin{proof}
Functional calculus gives
\(D_M^r=\exp((\log R_{\rm act})B_M^r/2)\), with strictly
positive spectrum. Its inverse is the exponential of the negative exponent.
For every \(v\), the transformed quadratic form equals
\((D_M^rv)^*M_M^{(0)}(D_M^rv)\). The invertible change of variable
preserves the maximal dimensions of positive, negative
and null subspaces; these determine inertia and rank. In a common
diagonal basis, the two factors jointly contribute
\(R_{\rm act}^{\kappa_\gamma^r}\). Dividing two values cancels
the common dimensional section and uses the multiplicative law of the
character.
\end{proof}

The distinction between transporting a form and modifying an
operator with respect to a fixed metric can be checked without approximations.
For
\[
 M_0=\begin{pmatrix}1&0\\0&4\end{pmatrix},\qquad
 D=\begin{pmatrix}2&0\\0&1/2\end{pmatrix},\qquad
 M'=D^*M_0D=\begin{pmatrix}4&0\\0&1\end{pmatrix},
\]
the determinant of \(D\) is one and the values associated with the two
routes are interchanged. Even the unordered spectrum changes if
\(M_0=I\): it changes from \(\{1,1\}\) to \(\{4,1/4\}\).
Both checks preserve positivity, rank and inertia.

\begin{proposition}[Spectral invariance with a transported metric]
Let \(G_0>0\) be a metric and let \(D\) be invertible. If
\[
 M'=D^*M_0D,\qquad G'=D^*G_0D,
\]
the pairs \((M',G')\) and \((M_0,G_0)\) have the same generalized
spectrum, with the same multiplicities.
\end{proposition}
\begin{proof}
One has
\[
 \det(M'-\lambda G')=|\det D|^2\det(M_0-\lambda G_0),
 \qquad (G')^{-1}M'=D^{-1}G_0^{-1}M_0D.
\]
The first equality preserves the roots and their multiplicities; the
second exhibits a similarity of the associated operators.
\end{proof}

The positive diagonal flow induced by \(K\) therefore has two precise
mathematical realizations on a material form. When it transports
the form and its metric simultaneously, it describes the same
object in another chart. When it acts on the form while preserving the route
metric, it produces an operator collection that may change the relative
masses. Its physical identification is determined by the route reader and
the declared action.

\subsection{Compatibility of the dodecaphasic flow with material return}

Let
\[
 A_K=(\log\rho_A)\operatorname{diag}(\widehat u_{K,1},\ldots,
       \widehat u_{K,12}),\quad
 \widehat u_K=u_K/\|u_K\|,\quad \rho_A=\pi/\varphi^2,
\]
so that \(g_K(t)=e^{tA_K}\). The material return collection
on the multisection \(M\) has generator
\[
 C_M^r=\tfrac12(\log R_{\rm act})B_M^r,\qquad
 S_M^r(t)=e^{tC_M^r}.
\]

\begin{proposition}[Transport of congruence by an intertwining reader]
A declared linear map
\(J_M:\mathbb C^{12}\to\mathscr H_M\) satisfying
\begin{equation}
 C_M^rJ_M=J_MA_K
 \label{eq:mass-k-intertwiner}
\end{equation}
satisfies
\[
 S_M^r(t)J_M=J_Mg_K(t),
\]
and, for \(M_M(t)=S_M^r(t)M_M^{(0)}S_M^r(t)\),
\begin{equation}
 J_M^*M_M(t)J_M
 =g_K(t)^*\bigl(J_M^*M_M^{(0)}J_M\bigr)g_K(t).
 \label{eq:mass-k-pullback}
\end{equation}
\end{proposition}
\begin{proof}
Identity \eqref{eq:mass-k-intertwiner} is iterated to each power
and summed in the convergent exponential series. Taking adjoints and
substituting in the material form proves
\eqref{eq:mass-k-pullback}.
\end{proof}

The criterion is finite and admits a direct operator-level
stress test. In the declared diagonal bases, if \(a_j\) and
\(c_\gamma\) are the eigenvalues of \(A_K\) and \(C_M^r\),
respectively, one must have
\[
 (c_\gamma-a_j)(J_M)_{\gamma j}=0
 \qquad\text{for each }(\gamma,j).
\]
Each nonzero coefficient of the reader requires the corresponding equality
of eigenvalues. The recovered composition in
\eqref{eq:mass-k-alpha}--\eqref{eq:mass-k-electron} and the reversible chart
\eqref{eq:mass-k-event-chart} are effective results. The additional
identification of their two flows requires a reader \(J_M\) with property
\eqref{eq:mass-k-intertwiner}; this proposition determines exactly
what a realization supporting that identification must verify.

\subsection{Elliptic transport and regular Legendre transformation}

The variational realization of transport has an explicit antecedent
in the action ellipse. The angular pair already obtained determines
\(\eta=\operatorname{artanh}(C^*/A)\), in the chamber
\(A\ne0\), \(|C^*/A|<1\). The coordinates \((Q,P)\), both with
the dimension of the square root of action, realize the form
\[
 \mathcal I_\eta(Q,P)=\frac12\bigl(e^{-\eta}Q^2+e^\eta P^2\bigr),
 \qquad \lambda=P\,dQ,\quad \Omega=dQ\wedge dP.
\]
The transport between forms, with
\(M_\eta=\operatorname{diag}(e^{\eta/2},e^{-\eta/2})\) and
\(J_2=\left(\begin{smallmatrix}0&-1\\1&0\end{smallmatrix}\right)\),
is \(M_{\eta'}e^{-\vartheta J_2}M_\eta^{-1}\). An
interpolation \(\eta(t)\), \(\vartheta(t)\), with
\(\dot\vartheta=\omega(t)>0\), has the Hamiltonian recovered
from Article X:
\begin{equation}
 H_{\rm tr}(Q,P,t)=\omega(t)\mathcal I_{\eta(t)}(Q,P)
                 +\frac{\dot\eta(t)}2QP.
 \label{eq:mass-k-hamiltonian}
\end{equation}
The second term comes from \(\dot M_\eta M_\eta^{-1}\).
The equations
\[
 \dot Q=\omega e^\eta P+\frac{\dot\eta}2Q,
 \qquad
 \dot P=-\omega e^{-\eta}Q-\frac{\dot\eta}2P
\]
preserve \(\mathcal I_{\eta(t)}(Q(t),P(t))\): the explicit
variation of \(\eta\) and the contribution of the connection term
cancel. This result and the identity
\(P\dot Q-H_{\rm tr}=p\dot q-\omega(q^2+p^2)/2\), for
\(Q=e^{\eta/2}q\), \(P=e^{-\eta/2}p\), are proved
in the section on variational conservation in the work on the holographic
extreme and mean ratio.

\begin{proposition}[Regular Lagrangian of elliptic transport]
Suppose \(\eta\in C^2\), \(\omega\in C^1\) and \(\omega>0\).
The Legendre transformation with respect to \(P\) is invertible and
yields
\begin{equation}
 L_{\rm tr}(Q,\dot Q,t)=
 \frac{\bigl(\dot Q-\dot\eta Q/2\bigr)^2}{2\omega e^\eta}
 -\frac{\omega e^{-\eta}}2Q^2.
 \label{eq:mass-k-lagrangian}
\end{equation}
\end{proposition}
\begin{proof}
Write \(a=\dot\eta/2\), \(b=\omega e^\eta>0\) and
\(c=\omega e^{-\eta}\). Then
\(H_{\rm tr}=bP^2/2+aQP+cQ^2/2\), and
\[
 \dot Q=\partial_PH_{\rm tr}=bP+aQ,
 \qquad P=\frac{\dot Q-aQ}{b}.
\]
Substituting this expression in \(P\dot Q-H_{\rm tr}\), the
terms containing \(P\) combine into
\(P(\dot Q-aQ)-bP^2/2=(\dot Q-aQ)^2/(2b)\), which
proves \eqref{eq:mass-k-lagrangian}. Moreover,
\[
 \partial_{\dot Q}L_{\rm tr}=P,
 \qquad
 \partial_P^2H_{\rm tr}=\omega e^\eta>0,
 \qquad
 \partial_{\dot Q}^2L_{\rm tr}=\frac1{\omega e^\eta}>0.
\]
The transformation is regular and its Euler--Lagrange equations
are equivalent to the preceding Hamiltonian equations.
\end{proof}

This algebraic elimination joins the Hamiltonian already constructed to
its second-order Lagrangian form. Regularity follows from the
quadratic term of the elliptic realization. For a purely linear
cotangent lift \(H_K(q,p)=p^{\mathsf T}A_Kq\), the
equation \(\dot q=A_Kq\) is independent of \(p\), and the Hessian
\(\partial_p^2H_K\) vanishes. Its variational formulation uses the
first-order action \(\int p^{\mathsf T}(\dot q-A_Kq)\,dt\).
The corpus also contains its discrete antecedent
\(L_n=p_{n+1}^{\mathsf T}(q_{n+1}-U_nq_n)\), whose equations
yield the transport \(q_{n+1}=U_nq_n\) and the dual transport
\(p_{n+1}=U_n^{-\mathsf T}p_n\). Both formulations have
specified domains and operators. The temporal interpolation and the
action section are retained as part of the declared physical
realization.
